\scanpage{153}
3. \uline{Application à la caractérisation des \deleted{\unsure{morphismes}} \inserted{revêtements} étales \unsure{locaux|loc.}} \deleted{\unsure{galoisiens}} \unsure{Cf plus bas.}

\inserted{$I_n,\ \unsure{\mathfrak S_n|\mathfrak G_n|G_n}$}

\uline{Prop 3.1.} Soit $f:X\to Y$ un morphisme fini, \inserted{\unsure{$n$ un entier|ou un entier}} Conditions équivalentes

(i) $f$ est étale, et le rang de $X$ aux p\textsuperscript{ts} de $Y$ est partout égal à $n$

(ii) Il existe un $Y_1\unsure{-|\to}Y$ fid plat tel que $X\times_Y Y_1$ soit $Y_1$-isomorphe à $Y_1\times I_n$.

(ii bis) Kif-kif, avec $Y_1$ \deleted{étale} rev. étale galoisien (de groupe $\mathfrak S_n$ si on \unsure{veut|semble})

(iii) $X$ est isomorphe à un $P/\deleted{\unsure{G}}H$, où $P$ est un rev. étale galoisien de $Y$, de groupe $G$, et $H$ un sous-groupe de $G$.

(iii bis) Kif kif, avec $G=\mathfrak S_n$, $H=\mathfrak S_{n-1}$, \inserted{(\unsure{voie|voir} d'indice $n$)} (stabilisateur d'un pt de $I_n$).

% slanted
\margin{(i) $\Rightarrow$ (iii bis) théorie de Galois,\\ (iii) $\Rightarrow$ (ii bis) immédiat car les \uline{foncteurs} des \uline{\unsure{associés}} \uline{est compatible} avec \uline{images} \uline{inverses}\\ (ii) $\Rightarrow$ (i) vu par la descente \unsure{$\{$}}

(i) $\Longrightarrow$ (iii bis) (La démonstration qui \deleted{suit} \inserted{ainsi que le résultat} \deleted{est} \unsure{due|dues} à \inserted{J.P.} Serre). On fait opérer $\mathfrak S_n$ sur $X\times\ldots\times X$ (les facteurs, produits sur $Y$). façon évidente. L'ensemble des pts d'inertie est \deleted{fermé et ouvert} \inserted{réunion de composantes} connexes (car si $s\in\mathfrak S_n$ est dans le groupe d'inertie d'un pt $\alpha$, il opère trivial\textsuperscript{t} sur

% archive sheet number: 152

